\clearpage\begingroup\SourceLeaf{77a}\NoSourceHead
\begin{center}\Large\textls[160]{BRIEFWECHSEL.}\end{center}
\bigskip\begin{center}[1.]\end{center}
\begin{center}\textsc{Gauss} an \textsc{Gerling}.\quad Göttingen, 14. Februar 1832.\end{center}
\noindent . . . . . . . . . . . . . . . . . . . .\par
\indent Ich habe mich in der letzten Zeit etwas mit dem Magnetismus überhaupt
beschäftigt, namentlich auch die Intensität des Erdmagnetismus auf eine absolute,
klar verständliche Einheit zu bringen gesucht. Ich finde, dass sie
immer die Form hat \(\dfrac{\sqrt g}{r}\), wo \(g\) eine Gewicht und \(r\) eine Linie ist. Nach
einigen Versuchen ist in Göttingen, wenn \(r\) ein Zoll ist, \(g\) nur \textls[120]{wenige
Milligramme} gross. Die Zeit ist heute zu kurz, mich weiter darüber zu
erklären, zumahl da meine Rechnungen noch nicht reif sind. . . . . . . . . .
\clearpage\endgroup
